\begin{frame}[shrink]
	\frametitle{O backward é a mesma matriz, transposta}
	\small
	\setlength{\abovedisplayskip}{2pt}
	\setlength{\belowdisplayskip}{2pt}
	\begin{itemize}
		\setlength{\itemsep}{0pt}
		\item No forward, $W_2$ levou 2 valores (os ocultos) para 1 (a saída). O gradiente precisa fazer o caminho \textbf{inverso}: 1 volta para 2. A matriz que faz isso é a \textbf{mesma}, só que transposta: o que era linha vira coluna. \newline \newline \textbf{O erro volta pelos próprios pesos do forward}, não por pesos novos.\newline 
		\item $W \leftarrow W - \eta \, \partial L/\partial W$\newline
	\end{itemize}
	
	\begin{equation*}
		\frac{\partial L}{\partial y} = W^{\top} \frac{\partial L}{\partial z}
		\qquad\qquad
		\frac{\partial L}{\partial W} = \frac{\partial L}{\partial z} \otimes \text{entrada}
	\end{equation*}
	\par A regra da cadeia diluída nessas matrizes é o produto das derivadas locais de cada elo do caminho $x_1 \rightarrow H_1 \rightarrow O \rightarrow L$:
	\begin{equation*}
		\frac{\partial L}{\partial w_{11}} =
		\underbrace{(-0{,}3173)}_{\partial L/\partial y_O} \!\cdot\!
		\underbrace{0{,}2432}_{\sigma'(z_O)} \!\cdot\!
		\underbrace{1{,}2}_{w_{H_1 O}} \!\cdot\!
		\underbrace{0{,}2314}_{\sigma'(z_{H_1})} \!\cdot\!
		\underbrace{0{,}9}_{x_1}
		= -0{,}01928
	\end{equation*}
\end{frame}
